% ============================================================================
% Chapter 16: Unified Picture
% ============================================================================
% Source: Synthesis of Chapters 1-15
% Status: FILLED
% Contamination check: PASSED (no banned terms)
% Contamination guard: Layer A uses EDC-native vocabulary only.
% ============================================================================

\chapter{Unified Picture}
\label{ch:unified}

% ----------------------------------------------------------------------------
\section*{Abstract}
% ----------------------------------------------------------------------------

This chapter synthesizes the derivation chain developed throughout Book IV.
From a single geometric source---the brane tension $\sigma$ inherited from
Book I---combined with 5D topology and the M$_6$ coordination lattice, we
derive: (1) junction stability and metastable transition lifetimes,
(2) pinning energies for small clusters, (3) closed-4 binding budgets,
(4) barrier-limited release systematics, and (5) coordination frustration
corrections that restore predictivity in the high-coordination regime.
Layer A maintains strict EDC-native vocabulary throughout; all calibrated
parameters are quarantined in Appendix~Q; and unresolved derivation gaps
are flagged as [OPEN]. This chapter introduces no new physics---it only
consolidates, cross-links, and assigns epistemic status to existing results.

% ============================================================================
\section{Ontology Recap}
\label{sec:unified:ontology}
% ============================================================================

Table~\ref{tab:unified:glossary} defines the core objects used throughout
Book IV in EDC-native terminology.

\begin{table}[h]
\centering
\caption{EDC-native ontology for Book IV.}
\label{tab:unified:glossary}
\begin{tabular}{lp{9cm}}
    \toprule
    Term & Definition \\
    \midrule
    Anchor junction &
    Z$_6$-symmetric Y-junction defect; topological ground state;
    stable against transition (Ch.~\ref{ch:anchor}) \\
    Metastable junction &
    Z$_3$-symmetric branch state; excited configuration in double-well
    potential; transitions to anchor via instanton tunneling (Ch.~\ref{ch:metastable}) \\
    Closed-4 unit &
    Minimal closed topological object with exceptional stability;
    4-constituent pinning cluster (Ch.~\ref{ch:helium4}) \\
    Pinning cluster &
    M$_6$-coordinated junction network; binding via topological
    localization (Ch.~\ref{ch:M6}) \\
    Coordination function $n(A)$ &
    Effective coordination number: $n(A) = p \cdot A^{1/3}$ (Ch.~\ref{ch:frustrated}) \\
    Allowed set $S$ &
    $\{2^a \times 3^b : a, b \geq 0\}$; coordination numbers compatible
    with Z$_6$ symmetry (Ch.~\ref{ch:frustrated}) \\
    Distance-to-allowed $d(n)$ &
    $\min_{k \in S}|n(A) - k|$; coordination frustration metric
    (Ch.~\ref{ch:frustrated}) \\
    Baseline lane &
    Empirical scaling law for barrier-limited release times [BL]
    (Ch.~\ref{ch:barrier_release}) \\
    Corrected lane &
    Baseline + frustration correction $g \cdot d(n)$
    (Ch.~\ref{ch:high_coord}) \\
    \bottomrule
\end{tabular}
\end{table}

\subsection{Inherited Parameters from Book I}

Book IV uses four parameters from Book I:
\begin{itemize}
    \item $\sigma$: Brane tension [I]
    \item $\delta$: 5D length scale [I]
    \item $L_0$: Characteristic brane separation [I]
    \item $T_*$: Temperature scale [I]
\end{itemize}

These are treated as inputs; their derivation belongs to Book I.

% ============================================================================
\section{The Derivation Tree}
\label{sec:unified:tree}
% ============================================================================

Table~\ref{tab:unified:tree} presents the complete derivation ledger for
Book IV, tracing how inputs propagate to observables.

\begin{table}[htbp]
\centering
\caption{Derivation Tree Ledger. Columns: Node identifier, input dependencies,
output result, source chapter, epistemic tags, and whether [Cal] is required.}
\label{tab:unified:tree}
\small
\begin{tabular}{lp{2.5cm}p{3.5cm}ccc}
    \toprule
    Node & Input(s) & Output & Ch. & Tags & Cal? \\
    \midrule
    A & 5D geometry & Z$_6$ anchor junction (Steiner minimum) &
        \ref{ch:anchor} & [Der] & N \\
    B & Z$_6$ structure & Z$_3$ subgroup, double-well $V(q)$ &
        \ref{ch:metastable} & [P] & N \\
    C & Node B, $V_B$ & Barrier height conjecture &
        \ref{ch:metastable} & [P] & N \\
    D & 5D action, Node C & Instanton: $\tau = A(\hbar/\omega_0)\exp(S_E/\hbar)$ &
        \ref{ch:instanton} & [Der] & N \\
    E & $\pi_1(S^1) = \mathbb{Z}$ & $\kappa = 2\pi$ (winding number) &
        \ref{ch:kappa} & [Dc] & N \\
    F & 5D geometry & $L_0/\delta = \pi^2$ &
        \ref{ch:L0delta} & [P] & N \\
    G & Nodes D,E,F & $S_E/\hbar = 2\pi^3$; $\tau_n \approx 880$ s &
        \ref{ch:tau_n} & [Dc]+[P] & Y \\
    H & $\sigma$, M$_6$ & Pinning constant $K$ &
        \ref{ch:sigma_K} & [Der] & N \\
    I & Node H & $B_2 = N_{\text{bonds}} K$; $N_{\text{bonds}} = 3$ &
        \ref{ch:deuterium} & [P] & N \\
    J & Node H, geometry & Closed-4 budget: 4 terms &
        \ref{ch:helium4} & [Der]+[P] & N \\
    K & Nodes I,J & Composition principle, emission bias &
        \ref{ch:light_clusters} & [Der]+[P] & N \\
    L & Empirical data & Baseline lane: $\log t = aX + b$ &
        \ref{ch:barrier_release} & [BL] & Y \\
    M & Z$_6$ symmetry & Allowed set $S = \{2^a 3^b\}$ &
        \ref{ch:frustrated} & [Der] & N \\
    N & Node M & Distance $d(n)$; forbidden zone &
        \ref{ch:frustrated} & [Der] & N \\
    O & Nodes L,N & Corrected lane: $\Delta = g \cdot d(n)$ &
        \ref{ch:frustrated} & [P]+[Cal] & Y \\
    P & Node O & High-coordination predictions &
        \ref{ch:high_coord} & [Dc] & Y \\
    \bottomrule
\end{tabular}

\vspace{2mm}
\footnotesize
\textbf{Legend:}
[Der] = derived from 5D; [Dc] = derived with constraints;
[P] = postulated/conjectured; [BL] = baseline empirical;
[Cal] = requires calibrated coefficients (Appendix~Q).
\end{table}

\subsection{Chain Segments}

\textbf{Segment 1: Junction Ground State (Ch.~\ref{ch:anchor})}
\begin{itemize}
    \item The Z$_6$-symmetric Y-junction emerges as the Steiner minimum
          of the 5D brane intersection problem [Der].
    \item This establishes absolute stability: no lower-energy configuration
          exists.
\end{itemize}

\textbf{Segment 2: Metastable Branch (Ch.~\ref{ch:metastable})}
\begin{itemize}
    \item Z$_3 \subset$ Z$_6$ identifies a higher-energy branch state [Der].
    \item The reaction coordinate $q \in [0,1]$ and double-well potential
          $V(q)$ are postulated [P].
    \item Barrier height $V_B$ conjecture connects to 5D geometry [P].
\end{itemize}

\textbf{Segment 3: Instanton Chain (Ch.~\ref{ch:instanton}--\ref{ch:tau_n})}
\begin{itemize}
    \item Instanton formula: $\tau = A(\hbar/\omega_0)\exp(S_E/\hbar)$ [Der]
    \item Topological factor: $\kappa = 2\pi$ from $\pi_1(S^1)$ [Dc]
    \item Scale ratio: $L_0/\delta = \pi^2$ [P] + [OPEN]
    \item Assembly: $S_E/\hbar = \kappa \cdot (L_0/\delta) = 2\pi^3$ [Dc]
    \item Prefactor $A$, $\omega_0$: partially [Cal], partially [OPEN]
    \item Result: $\tau_n \approx 880$ s [Dc]+[P]+[Cal]
\end{itemize}

\textbf{Segment 4: Pinning and Binding (Ch.~\ref{ch:deuterium}--\ref{ch:light_clusters})}
\begin{itemize}
    \item Pinning constant $K$ from $\sigma$ and contact geometry [Der]
    \item $A = 2$: $B_2 = 3K$ with $N_{\text{bonds}} = 3$ [P]
    \item $A = 4$: Four-term budget (localization + pinning + surface +
          closure) [Der]+[P]+[OPEN]
    \item Composition principle for $A \leq 10$ [Der]
    \item Closed-4 emission bias [P]
\end{itemize}

\textbf{Segment 5: Release Systematics (Ch.~\ref{ch:barrier_release}--\ref{ch:high_coord})}
\begin{itemize}
    \item Baseline lane from empirical scaling [BL]
    \item Allowed set $S$ from Z$_6$ symmetry [Der]
    \item Distance $d(n)$ and forbidden zone [37,47] [Der]
    \item Correction: $\Delta = g \cdot d(n)$ with $g < 0$ [P]+[Cal]
    \item High-coordination predictions: 7$\times$ error reduction [Dc]
\end{itemize}

% ============================================================================
\section{Epistemic Ledger Summary}
\label{sec:unified:epistemic}
% ============================================================================

Table~\ref{tab:unified:epistemic} summarizes the epistemic status of all
key quantities in Book IV.

\begin{table}[htbp]
\centering
\caption{Epistemic Status Summary. [I] = input from Book I;
[Der] = derived; [Dc] = derived with constraints; [P] = postulated;
[Cal] = calibrated (Appendix~Q); [OPEN] = derivation gap.}
\label{tab:unified:epistemic}
\small
\begin{tabular}{lccl}
    \toprule
    Quantity & Tag(s) & Cal? & Notes \\
    \midrule
    \multicolumn{4}{l}{\textit{Book I imports}} \\
    $\sigma$ (brane tension) & [I] & N & Fundamental input \\
    $\delta$ (5D scale) & [I] & N & Fundamental input \\
    $L_0$ (separation) & [I] & N & Fundamental input \\
    $T_*$ (temperature) & [I] & N & Fundamental input \\
    \midrule
    \multicolumn{4}{l}{\textit{Junction structure}} \\
    Z$_6$ ground state & [Der] & N & Steiner minimum \\
    Z$_3$ metastable & [Der] & N & Subgroup structure \\
    Double-well $V(q)$ & [P] & N & Conjecture \\
    Barrier $V_B$ & [P] & N & Conjecture \\
    \midrule
    \multicolumn{4}{l}{\textit{Instanton parameters}} \\
    $\kappa = 2\pi$ & [Dc] & N & From homotopy \\
    $L_0/\delta = \pi^2$ & [P]+[OPEN] & N & Geometry conjecture \\
    $\omega_0$ (frequency) & [P]+[OPEN] & Y & Prefactor block \\
    $A$ (attempt rate) & [P]+[OPEN] & Y & Prefactor block \\
    $S_E/\hbar = 2\pi^3$ & [Dc] & N & Assembled \\
    $\tau_n \approx 880$ s & [Dc]+[P] & Y & Final result \\
    \midrule
    \multicolumn{4}{l}{\textit{Pinning and binding}} \\
    $K$ (pinning constant) & [Der] & N & From $\sigma$, geometry \\
    $N_{\text{bonds}}(A{=}2) = 3$ & [P] & N & Bond count \\
    Closed-4 localization & [Der] & N & 21 MeV term \\
    Closed-4 pinning & [Der] & N & 5 MeV term \\
    Closed-4 surface & [P]+[OPEN] & N & 2 MeV term \\
    Closed-4 closure & [P]+[OPEN] & N & 2 MeV term \\
    \midrule
    \multicolumn{4}{l}{\textit{Release systematics}} \\
    Baseline $a$, $b$ & [BL] & Y & Appendix~Q \\
    Prefactor $p$ in $n(A)$ & [Cal] & Y & Appendix~Q \\
    Allowed set $S$ & [Der] & N & From Z$_6$ \\
    Distance $d(n)$ & [Der] & N & Algorithm \\
    Coefficient $g$ & [Cal] & Y & Appendix~Q \\
    Sign $g < 0$ & [P] & N & Frustration mechanism \\
    \bottomrule
\end{tabular}
\end{table}

\subsection{Promotion Paths}

To promote [P] quantities to [Der], the following derivation gaps must
be closed:

\begin{enumerate}
    \item \textbf{$L_0/\delta = \pi^2$:} Derive from 5D action with
          explicit boundary conditions.
    \item \textbf{Prefactor block ($A$, $\omega_0$):} Derive from
          fluctuation determinant around the instanton.
    \item \textbf{$N_{\text{bonds}} = 3$:} Prove from junction topology
          for $A = 2$ pinning cluster.
    \item \textbf{Closed-4 surface and closure terms:} Derive from
          variational principle on M$_6$ lattice.
    \item \textbf{Sign of $g$:} Derive from enhanced preformation or
          barrier reduction mechanism.
    \item \textbf{Prefactor $p$:} Derive $n(A) = p A^{1/3}$ from 5D
          dynamics rather than calibration.
\end{enumerate}

% ============================================================================
\section{What the Theory Explains}
\label{sec:unified:explains}
% ============================================================================

\subsection{Explained (Layer A Derivations)}

\begin{itemize}
    \item \textbf{Anchor junction stability:} Z$_6$ Steiner minimum
          provides absolute topological stability [Der].
    \item \textbf{Metastable junction existence:} Z$_3$ subgroup identifies
          the excited branch state [Der].
    \item \textbf{Pinning mechanism:} Brane tension $\sigma$ generates
          binding via topological localization [Der].
    \item \textbf{Allowed coordination set:} $S = \{2^a \times 3^b\}$
          follows from Z$_6$ symmetry [Der].
    \item \textbf{Forbidden zone:} Gap [37,47] is a direct consequence
          of allowed-set structure [Der].
\end{itemize}

\subsection{Partially Explained (Derivation + Postulate)}

\begin{itemize}
    \item \textbf{Metastable lifetime scale:} $\tau_n \approx 880$ s
          from instanton chain, but prefactor block is [OPEN].
    \item \textbf{$A = 2$ binding:} $B_2 = 3K$ matches observations, but
          $N_{\text{bonds}} = 3$ is [P].
    \item \textbf{Closed-4 binding budget:} 4-term decomposition, but
          surface and closure terms are [P]+[OPEN].
    \item \textbf{High-coordination predictivity:} 7$\times$ improvement
          established [Dc], but $g$ mechanism is [P].
\end{itemize}

\subsection{Open (Not Yet Derived)}

\begin{itemize}
    \item Triangular ($A = 3$) binding energy derivation
    \item Boundary reconfiguration dynamics
    \item Alternative emission channel competition
    \item First-principles derivation of $p$, $g$, and prefactor block
\end{itemize}

% ============================================================================
\section{Open Problems Register}
\label{sec:unified:open}
% ============================================================================

\begin{edcopenbox}[title=Consolidated Open Problems]
\textbf{5D Action Derivations:}
\begin{enumerate}
    \item Derive $L_0/\delta = \pi^2$ from bulk + brane + GHY + Israel
          junction conditions (Ch.~\ref{ch:L0delta}).
    \item Derive double-well potential $V(q)$ and barrier height $V_B$
          from 5D dynamics (Ch.~\ref{ch:metastable}).
    \item Derive coordination function $n(A) = p A^{1/3}$ with $p$
          determined by 5D geometry (Ch.~\ref{ch:frustrated}).
\end{enumerate}

\textbf{Prefactor Derivations:}
\begin{enumerate}
    \setcounter{enumi}{3}
    \item Derive attempt frequency $\omega_0$ from brane fluctuation
          spectrum (Ch.~\ref{ch:instanton}).
    \item Derive prefactor $A$ from fluctuation determinant around
          instanton (Ch.~\ref{ch:tau_n}).
\end{enumerate}

\textbf{Binding Structure:}
\begin{enumerate}
    \setcounter{enumi}{5}
    \item Prove $N_{\text{bonds}} = 3$ for $A = 2$ from junction topology
          (Ch.~\ref{ch:deuterium}).
    \item Derive closed-4 surface and closure terms from variational
          principle (Ch.~\ref{ch:helium4}).
    \item Complete triangular ($A = 3$) binding derivation
          (Ch.~\ref{ch:light_clusters}).
\end{enumerate}

\textbf{Release Mechanism:}
\begin{enumerate}
    \setcounter{enumi}{8}
    \item Derive frustration coefficient $g$ and its sign from
          preformation or barrier-reduction mechanism (Ch.~\ref{ch:frustrated}).
    \item Explain boundary reconfiguration dynamics not captured by
          scalar $d(n)$ (Ch.~\ref{ch:high_coord}).
    \item Characterize alternative emission channel competition
          (Ch.~\ref{ch:high_coord}).
\end{enumerate}

\textbf{Cross-Book:}
\begin{enumerate}
    \setcounter{enumi}{11}
    \item Complete Book I $\to$ Book IV parameter flow with explicit
          numerical chain from $\sigma$ to all observables.
\end{enumerate}
\end{edcopenbox}

% ============================================================================
\section{Reader Path}
\label{sec:unified:path}
% ============================================================================

Table~\ref{tab:unified:path} suggests learning paths through Book IV
depending on reader goals.

\begin{table}[h]
\centering
\caption{Suggested reading paths.}
\label{tab:unified:path}
\begin{tabular}{lp{8cm}}
    \toprule
    Path & Chapters \\
    \midrule
    \textbf{Results only} &
    Ch.~\ref{ch:anchor} (stability),
    Ch.~\ref{ch:tau_n} (lifetime),
    Ch.~\ref{ch:deuterium}--\ref{ch:helium4} (binding),
    Ch.~\ref{ch:high_coord} (predictions),
    Ch.~\ref{ch:unified} (synthesis) \\
    \textbf{Full derivation} &
    Ch.~\ref{ch:anchor}, \ref{ch:metastable},
    Ch.~\ref{ch:instanton}--\ref{ch:tau_n},
    Ch.~\ref{ch:deuterium}--\ref{ch:light_clusters},
    Ch.~\ref{ch:barrier_release}--\ref{ch:high_coord},
    Ch.~\ref{ch:unified} \\
    \textbf{Calibration audit} &
    Appendix~Q (all sections),
    Ch.~\ref{ch:repro} (reproducibility) \\
    \bottomrule
\end{tabular}
\end{table}

% ============================================================================
% Observer-side projection
% ============================================================================

\begin{observerbox}
Observer-side projection note: quoted terms are measurement labels for the
3D/4D projection of the 5D objects defined in this book; they do not
imply any conventional mechanism.

\begin{itemize}
    \item \AnchorJunction{} $\leftrightarrow$ ``proton'' (projection label)
    \item \MetastableJunction{} $\leftrightarrow$ ``neutron'' (projection label)
    \item \ClusterState{} $\leftrightarrow$ bound system (projection label)
    \item \ClosedFour{} $\leftrightarrow$ released unit (projection label)
\end{itemize}

What the observer records: the unified derivation chain maps 5D parameters
to projection-level observables---masses, lifetimes, binding energies,
and release times---providing consistency checks against measurement.
\end{observerbox}

% ============================================================================
\section{Summary and Forward Links}
\label{sec:unified:summary}
% ============================================================================

\begin{resultbox}
    \textbf{Book IV Synthesis:}
    \begin{itemize}
        \item From brane tension $\sigma$ (Book I) + 5D topology + M$_6$
              lattice, we derive junction stability, metastable lifetimes,
              pinning energies, and release systematics.
        \item The derivation tree (Table~\ref{tab:unified:tree}) traces
              16 nodes from inputs to observables.
        \item The epistemic ledger (Table~\ref{tab:unified:epistemic})
              distinguishes [Der], [Dc], [P], [Cal], and [OPEN] claims.
        \item 12 open problems are consolidated for future work.
    \end{itemize}
\end{resultbox}

\bigskip
\noindent
\textbf{Forward references:}
\begin{itemize}
    \item Chapter~\ref{ch:repro}: Reproducibility procedures and
          verification protocols.
    \item Appendix~Q: All calibration details, datasets, and fitting
          procedures are quarantined here.
    \item Appendix~X: Optional mapping to conventional nomenclature
          (non-binding, for reader orientation only).
\end{itemize}

\bigskip
\noindent
\textbf{Reader Contract:}
Layer A of this book uses EDC-native vocabulary exclusively. All
conventional mappings (junction $\leftrightarrow$ conventional labels,
etc.) are optional and quarantined in Appendix~X. All calibrated
parameters ([Cal]) are documented in Appendix~Q with anti-tuning
guardrails.

